\chapter{GIỚI THIỆU TỔNG QUAN}
\label{chap:intro}

\section{Bối cảnh: Sự phát triển của Tương tác Cộng tác Người - Máy (HRC)}
Trong những thập kỷ gần đây, ngành công nghiệp tự động hóa đã chứng kiến những bước tiến đáng kể. Chức năng của robot không chỉ giới hạn ở việc hoạt động độc lập bên trong các lồng kính an toàn tại các dây chuyền sản xuất hàng loạt, mà đang dần chuyển dịch sang mô hình Tương tác cộng tác Người - Máy (Human-Robot Collaboration - HRC) \cite{wang2024genealogy}. Mục tiêu của sự dịch chuyển này là kết hợp sự bền bỉ, tính chính xác của máy móc với khả năng nhận thức, ra quyết định linh hoạt của con người trước những tình huống phi cấu trúc.

Sự cộng tác này đặc biệt cần thiết trong các môi trường làm việc đặc thù như sản xuất linh kiện phức tạp, y tế, xử lý vật liệu độc hại, hay không gian vũ trụ. Tại các môi trường này, việc yêu cầu con người trực tiếp có mặt tại hiện trường mang lại nhiều rủi ro về an toàn \cite{autolab2021telesurgery}. Mặc dù các công nghệ Trí tuệ nhân tạo (AI) và Học máy đang phát triển, các hệ thống robot hiện tại vẫn gặp khó khăn trong việc tự chủ giải quyết các tình huống bất ngờ nằm ngoài dữ liệu huấn luyện. Do đó, phương pháp Điều khiển từ xa (Teleoperation) tiếp tục là một giải pháp thiết thực, cho phép chuyên gia con người giám sát và trực tiếp vận hành hệ thống robot từ một khoảng cách an toàn \cite{darvish2023teleoperation}.

\section{Ứng dụng Thực tế ảo (VR) trong Điều khiển từ xa}
Sự thành công của một hệ thống điều khiển từ xa phụ thuộc phần lớn vào Giao diện tương tác Người - Máy (Human-Machine Interface). Trong các hệ thống điều khiển công nghiệp truyền thống, người vận hành thường phải sử dụng màn hình 2D, bàn phím, hoặc bộ điều khiển cầm tay (Teach Pendant) để tương tác với robot \cite{hetrick2020comparing}. Quá trình này đòi hỏi người dùng phải liên tục quan sát không gian qua camera 2D, đối chiếu với mô hình 3D trong tư duy và thực hiện các thao tác cơ học. Điều này vô hình trung tạo ra tải nhận thức (cognitive load) lớn, dễ gây mệt mỏi và làm giảm độ chính xác khi thực hiện các tác vụ phức tạp \cite{chen2023comparing}.

Để khắc phục những hạn chế của giao diện 2D, công nghệ Thực tế ảo (Virtual Reality - VR) đã được nghiên cứu và ứng dụng như một giải pháp thay thế hiệu quả. Không gian VR cung cấp khả năng hiển thị chiều sâu tự nhiên, cho phép người dùng quan sát môi trường làm việc của robot dưới dạng 3D thông qua kính hiển thị đội đầu (HMD) \cite{arclab2025beavr}. Khi kết hợp với công nghệ Bản sao số (Digital Twin), người vận hành có thể tương tác trực quan với một mô hình ảo của robot được đồng bộ trạng thái với robot thực tế. Phương pháp này hỗ trợ cải thiện nhận thức không gian (spatial awareness) và giúp các thao tác vận hành trở nên tự nhiên hơn \cite{torrejon2026teleoperation}. 

\section{Giới hạn công nghệ và Giải pháp tiếp cận}
Mặc dù việc tích hợp VR mang lại nhiều hứa hẹn cho hệ thống điều khiển từ xa, quá trình triển khai thực tế trên các dòng robot công nghiệp phổ thông (như Denso VS-6577) vẫn tồn tại những rào cản nhất định về cả nền tảng phần mềm lẫn khả năng tương thích của phần cứng.

\subsection{Sự hạn chế của nền tảng ROS 1 trong ứng dụng thời gian thực}
Robot Operating System (ROS 1) đã từ lâu đóng vai trò là một nền tảng tiêu chuẩn trong nghiên cứu robot. Tuy nhiên, kiến trúc của ROS 1 bộc lộ một số hạn chế khi triển khai các hệ thống yêu cầu tính thời gian thực (real-time). ROS 1 sử dụng kiến trúc tập trung với một Master Node duy nhất và giao thức truyền thông TCPROS \cite{macenski2022robot}. Trong hệ thống điều khiển qua VR, dữ liệu truyền tải bao gồm cả luồng hình ảnh camera (băng thông lớn) và tín hiệu điều khiển động học (kích thước nhỏ nhưng cần tần số cao). Việc truyền tải đồng thời qua kiến trúc của ROS 1 dễ dẫn đến tình trạng nghẽn mạng cục bộ hoặc trễ gói tin (packet loss), từ đó làm giảm độ mượt mà của thao tác điều khiển \cite{kawshan2026digital, casini2025survey}. 

Để giải quyết vấn đề này, sự dịch chuyển sang ROS 2 là một yêu cầu tự nhiên. Nhờ sử dụng giao thức Data Distribution Service (DDS), ROS 2 cung cấp cơ chế mạng ngang hàng phân tán và các cấu hình Chất lượng dịch vụ (QoS) linh hoạt. Điều này cho phép hệ thống phân luồng và ưu tiên các gói tin điều khiển quan trọng, góp phần duy trì tính ổn định của hệ thống trong môi trường truyền thông đa luồng \cite{macenski2022robot, ali2025digital}.

\subsection{Bài toán tích hợp với thiết bị phần cứng đời cũ}
Hầu hết các nghiên cứu viễn thao hiện đại thường được thực hiện trên các nền tảng robot nghiên cứu thế hệ mới, vốn được hỗ trợ sẵn các giao thức điều khiển thời gian thực. Trong khi đó, ở môi trường công nghiệp, có rất nhiều robot thế hệ cũ (như bộ điều khiển RC5 của Denso) đang được vận hành. Các thiết bị này chủ yếu giao tiếp qua chuẩn UART truyền thống và được thiết kế để nhận tập lệnh lập trình sẵn (offline) thay vì nhận tín hiệu điều khiển liên tục (real-time servoing). Sự chênh lệch giữa tốc độ phát lệnh từ phần mềm hiện đại và khả năng đáp ứng cơ khí của phần cứng cũ dễ dẫn đến tình trạng mất đồng bộ, gây ra hiện tượng giật cục trong chuyển động. Đây là một bài toán thực tiễn cần có giải pháp tích hợp phần mềm trung gian để duy trì sự tương thích.

\section{Mục tiêu và Giải pháp của đề tài}
Từ các phân tích trên, mục tiêu của luận văn là nghiên cứu và xây dựng một hệ thống phần mềm điều khiển cánh tay robot công nghiệp Denso VS-6577 từ xa thông qua kính thực tế ảo Meta Quest 3, sử dụng nền tảng ROS 2 làm cốt lõi giao tiếp. Đề tài tập trung vào việc tạo ra một chu trình kết nối liền mạch từ giao diện người dùng đến thiết bị phần cứng, thông qua các giải pháp cụ thể sau:

\begin{itemize}
    \item \textbf{Thiết kế kiến trúc phần mềm dựa trên ROS 2 và MoveIt 2:} Xây dựng hệ thống điều khiển tách biệt (decoupled) theo tiêu chuẩn Ros2-Control. Việc áp dụng ROS 2 hỗ trợ phân luồng dữ liệu hiệu quả hơn, đảm bảo tín hiệu điều khiển và dữ liệu hình ảnh từ không gian làm việc được truyền tải đến kính VR một cách ổn định, khắc phục giới hạn băng thông của hệ thống ROS 1.
    \item \textbf{Phát triển ứng dụng VR tương tác trực quan:} Xây dựng ứng dụng điều khiển trên thiết bị Meta Quest 3 với engine Unity 6. Ứng dụng cung cấp cho người dùng một giao diện Digital Twin đồng bộ trạng thái theo thời gian thực và hiển thị luồng video từ camera, giúp tăng cường khả năng quan sát và hỗ trợ người dùng đưa ra các thao tác di chuyển cánh tay linh hoạt.
    \item \textbf{Xây dựng lớp giao tiếp trung gian cho thiết bị công nghiệp:} Để giải quyết bài toán giao tiếp với bộ điều khiển RC5 đời cũ, đề tài phát triển một module Hardware Interface chuyên biệt. Module này ứng dụng kỹ thuật bộ đệm (Ring Buffer) và xử lý xu hướng tín hiệu để làm mượt các lệnh điều khiển được gửi xuống từ phần mềm. Mặc dù đây chỉ là các kỹ thuật xử lý tín hiệu tạm thời và không thể loại bỏ hoàn toàn độ trễ cơ khí vật lý, giải pháp này góp phần duy trì khả năng hoạt động ổn định của robot, hạn chế tối đa các lỗi giật cục do sự thiếu đồng bộ giữa phần mềm điều khiển và phần cứng.
\end{itemize}

\section{Bố cục của luận văn}
Luận văn được tổ chức thành 6 chương với nội dung chính như sau:
\begin{itemize}
    \item \textbf{Chương 1: Giới thiệu tổng quan.} Trình bày bối cảnh chung, giới thiệu vai trò của công nghệ VR, các hạn chế kỹ thuật hiện tại và giải pháp mà đề tài hướng đến.
    \item \textbf{Chương 2: Các nghiên cứu liên quan.} Lược khảo và đánh giá các công trình nghiên cứu hiện có về giao diện VR, bản sao số và tối ưu nền tảng middleware trong điều khiển robot từ xa.
    \item \textbf{Chương 3: Cơ sở lý thuyết.} Cung cấp kiến thức nền tảng về động học robot, kiến trúc ROS 2, tiêu chuẩn Ros2-Control và các công cụ phát triển phần mềm (Unity 6, Meta Quest 3).
    \item \textbf{Chương 4: Thiết kế và hiện thực hệ thống.} Trình bày chi tiết sơ đồ khối, cách thức xây dựng module giao tiếp UART, kiến trúc ROS 2 và cách phát triển ứng dụng điều khiển VR.
    \item \textbf{Chương 5: Kết quả hiện thực và đánh giá.} Đưa ra các minh chứng thực nghiệm về giao diện, đánh giá hiệu năng phần mềm và khả năng đáp ứng của hệ thống phần cứng.
    \item \textbf{Chương 6: Tổng kết và hướng phát triển.} Tóm tắt kết quả đạt được, thừa nhận những điểm hạn chế còn tồn tại và đề xuất hướng nghiên cứu tiếp theo.
\end{itemize}

% =========================================================================================

\chapter{CÁC NGHIÊN CỨU LIÊN QUAN}
\label{chap:related_works}

Việc ứng dụng Thực tế ảo (VR) vào quá trình viễn thao (Teleoperation) đã và đang nhận được nhiều sự quan tâm từ cộng đồng nghiên cứu robot. Chương này sẽ trình bày tổng quan các nghiên cứu liên quan, được phân loại theo ba nhóm khía cạnh chính: (1) Phương thức tương tác và giao diện điều khiển, (2) Khả năng cải thiện nhận thức tình huống qua Bản sao số, và (3) Vai trò của các nền tảng Middleware trong việc duy trì kết nối hệ thống. Việc phân tích này đóng vai trò làm cơ sở để định vị giải pháp kỹ thuật của luận văn trong bức tranh công nghệ hiện tại.

\section{Phương thức tương tác và giao diện điều khiển (Control Interfaces)}
Mục tiêu cơ bản của hệ thống viễn thao là thu hẹp khoảng cách không gian giữa người vận hành và thiết bị. Ở khía cạnh này, việc lựa chọn phương thức nhập liệu (input methods) đóng vai trò quan trọng đối với trải nghiệm người dùng. 

Nghiên cứu của Chen và cộng sự (2023) \cite{chen2023comparing} đã thực hiện đánh giá ba phương thức điều khiển trên cánh tay JAKA Minicobo: giao diện 2D (GUI), nhận diện cử chỉ tay (Hand Tracking), và bộ điều khiển VR (VR Controllers). Kết quả thực nghiệm cho thấy việc sử dụng VR Controllers giúp người dùng hoàn thành nhiệm vụ nhanh chóng và ổn định hơn so với hai phương pháp còn lại, phần lớn nhờ vào khả năng theo dõi chuyển động không gian 6 bậc tự do (6-DoF) đáng tin cậy. Dù cử chỉ tay mang lại cảm giác công thái học tự nhiên, chúng dễ bị nhiễu tín hiệu hoặc mất theo dõi trong một số góc khuất. Phù hợp với quan điểm này, hệ thống BEAVR của Viện MIT (2025) \cite{arclab2025beavr} cũng đã ứng dụng kính Meta Quest 3S và VR Controllers để điều khiển các thao tác phức tạp, mang lại hiệu quả tốt đối với người dùng phổ thông.

Gần đây, một số nghiên cứu còn mở rộng sang hướng tích hợp Phản hồi xúc giác (Haptic Feedback). Các nghiên cứu của Xu et al. (2025) \cite{xu2025immersive} và Papakonstantinou et al. (2025) \cite{papakonstantinou2025effect} cho thấy việc sử dụng thiết bị haptic giúp tăng độ chính xác trong các tác vụ thao tác tỉ mỉ. Tuy nhiên, việc áp dụng công nghệ này đòi hỏi cảm biến lực phức tạp và làm tăng đáng kể độ trễ truyền dẫn của hệ thống. Trong giới hạn phần cứng tiêu chuẩn của các cánh tay công nghiệp không trang bị cảm biến lực, đề tài hiện tại tập trung tối ưu hóa phản hồi thị giác (Visual Feedback) qua mô hình 3D và luồng video để bù đắp, đảm bảo tính khả thi khi triển khai thực tế.

\section{Cải thiện nhận thức tình huống thông qua Bản sao số (Digital Twin)}
Hạn chế của góc nhìn camera 2D tĩnh là thường tạo ra điểm mù khi làm việc trong không gian hẹp. Việc đồng bộ một Bản sao số trong môi trường VR giúp người dùng dễ dàng theo dõi toàn cảnh hệ thống.

Trong dự án \textit{Vicarios}, Naceri và cộng sự (2021) \cite{naceri2021vicarios} đã xây dựng một mô hình ảo của cánh tay UR5 sử dụng Unreal Engine và ROS. Đặc tính nổi bật của hệ thống là cho phép người dùng thay đổi góc nhìn (teleporting) để kiểm tra không gian làm việc từ nhiều hướng khác nhau. Gần đây hơn, Torrejón et al. (2026) \cite{torrejon2026teleoperation} đã tiến hành đồng bộ hóa Bản sao số kết hợp Đám mây điểm (Point Cloud) với mật độ dày đặc để tái tạo nguyên bản môi trường từ xa. Mặc dù đem lại cảm quan tốt, phương pháp truyền Point Cloud đòi hỏi phần cứng xử lý đồ họa cường độ cao và băng thông mạng lớn, đôi khi dẫn đến độ trễ hệ thống khá cao (lên tới 190ms).

Khác với định hướng mô phỏng môi trường nặng về dữ liệu, phương pháp của luận văn hiện tại sử dụng mô hình 3D kết xuất trực tiếp trên thiết bị (Meta Quest 3) và chỉ truyền tải dữ liệu trạng thái khớp (Joint States). Giải pháp này góp phần làm giảm tải cho đường truyền mạng, giúp ứng dụng VR duy trì tốc độ khung hình ổn định mà không đòi hỏi thiết bị máy tính đồ họa chuyên dụng đi kèm.

\section{Vai trò của Middleware và bài toán đồng bộ hệ thống}
Một khía cạnh thiết yếu trong nghiên cứu viễn thao là việc lựa chọn nền tảng giao tiếp (Middleware) phù hợp.

\begin{itemize}
    \item \textbf{Các hệ thống dựa trên ROS 1:} Nhiều dự án hiện nay vẫn sử dụng nền tảng ROS 1 \cite{naceri2021vicarios, chen2023comparing, kawshan2026digital}. Nghiên cứu của Kawshan và Peng (2026) \cite{kawshan2026digital} khi đánh giá kiến trúc tích hợp Unity-ROS 1 đã chỉ ra rằng giao thức mạng tập trung của ROS 1 thường sinh ra độ trễ tích lũy khi xử lý các tín hiệu yêu cầu tốc độ cao. 
    \item \textbf{Xu hướng dịch chuyển sang ROS 2:} Trong báo cáo tổng quan của Casini et al. (2025) \cite{casini2025survey}, ROS 2 với chuẩn giao tiếp DDS được đánh giá là cung cấp năng lực điều phối luồng dữ liệu ưu việt hơn, cho phép các hệ thống robot hoạt động với mức độ tin cậy thời gian thực. Ứng dụng điều này, Ali và cộng sự (2025) \cite{ali2025digital} đã triển khai ROS 2 để kết nối Digital Twin trong mô phỏng huấn luyện học tăng cường, đạt được mức độ trễ đồng bộ thấp. 
\end{itemize}

Dựa trên xu hướng này, luận văn tích hợp ROS 2 (Ros2-Control) làm nền tảng truyền thông cốt lõi giữa ứng dụng VR và khối xử lý trung tâm, nhằm tận dụng cơ chế phân luồng dữ liệu ổn định của DDS. 

\section{Nhận xét chung và sự định vị của luận văn}
Qua việc lược khảo các nghiên cứu trên, có thể thấy việc sử dụng VR và Bản sao số để điều khiển robot mang lại lợi ích rõ rệt trong việc cải thiện sự trực quan và hiệu năng thao tác. Các nghiên cứu hiện tại chủ yếu được thực hiện trên những nền tảng robot thế hệ mới với khả năng hỗ trợ lập trình điều khiển nhúng linh hoạt.

Đối với môi trường công nghiệp, việc tái sử dụng các hệ thống phần cứng đời cũ (chẳng hạn như bộ điều khiển RC5 trong luận văn này) thường đặt ra các bài toán thực tế về khả năng tương thích. Do đặc thù giao tiếp cơ bản (UART) và tốc độ đáp ứng cơ khí không được thiết kế cho việc điều khiển liên tục (real-time servoing), hệ thống dễ gặp phải sự sai lệch về mặt thời gian giữa luồng dữ liệu phát đi từ phần mềm và chuyển động của thiết bị vật lý. 

Luận văn hiện tại đóng vai trò là một dự án \textbf{tích hợp hệ thống mang tính thực tiễn}. Luận văn không tập trung vào việc tạo ra các thuật toán điều khiển mới hay thay thế hoàn toàn cấu trúc cơ khí, mà sử dụng công nghệ phần mềm hiện đại (ROS 2 và Unity VR) để cải thiện trải nghiệm vận hành của người dùng. Đồng thời, việc bổ sung các kỹ thuật lập trình đệm dữ liệu (buffer) ở cấp độ giao tiếp phần cứng được triển khai như một giải pháp hỗ trợ sự ổn định. Bằng cách thiết kế các module phân tách rõ ràng, luận văn kỳ vọng cung cấp một khuôn mẫu có khả năng điều chỉnh và ứng dụng cho việc nâng cấp giao diện điều khiển của các thiết bị công nghiệp truyền thống.
